\scanpage{29}

Le cas \deleted{\unsure{fidèle}} fidèle (resp. str. fidèle) est celui où ce sous-Module est un Module inversible, resp. un sous-fibré en droites. Dans ce dernier cas, il \deleted{\unsure{y}} y a un \deleted{fibré}\inserted{rev.} quadratique associé de $X$, qui peut s'interpréter comme l'intersection de $\unsure{\check P|\check{\mathbb P}}(\mathcal A)\subset\check P(\mathcal M)$ (où $\mathcal A\inserted{\supset\underline O\cdot1_{\mathcal M}}$ est le sous-Anneau de $\mathcal M$ image inverse de $\Delta$) avec la quadrique

\[
\mathcal Q\subset\check P(\mathcal M)
\]

définie par l'équation $\operatorname{Nr}_{\mathcal M/\unsure{k|h}}(x)=0$, où $\operatorname{Nr}$ désigne la norme réduite. (C'est tautologique, vu que ladite norme induit sur $\mathcal A$ celle de $\mathcal A$.)

Que devient le $\mathcal A$-Module inv \inserted{$E$}, i.e. le $\underline O_{X'}$-Module inversible $\Lambda$, de \unsure{la théorie précédente|le théorème précédent} ?

\unsure{Il faut|Il} s'attendre à ce qu'il ne soit plus défini que modulo remplacement de $E$ par $E\otimes_{\underline O_X}\Theta$, $\Theta$ inv. sur $X$, i.e. que $\Lambda$ soit défini mod changement en $\Lambda\otimes p^*(\Theta)$, où $p:X'\mathbin{\unsure{-|\longrightarrow}}X$ le morphisme structural. Mais si $\sigma$ désigne l'inv. can. de $X'/X$, alors $\Lambda\sigma{\unsure{(\Lambda)|({}^*\Lambda)|^*(\Lambda)}}^{-1}$ reste invariante, par changement de $\Lambda$ en $\Lambda\otimes p^*(\Theta)$. D'autre part,
% archive sheet number: 28
